\section{Uniform estimates and compactness}
\label{app:uniform}

The fitted velocity may change when the number of steps changes. We therefore
need estimates with the same constants for every velocity in the allowed
class. This is what ``uniform'' means below. An estimate for one fixed
velocity would not justify optimizing over the whole class. We derive these
estimates from the numerical recursion itself, without a Wasserstein tangent
expansion.\AIassistance

All norms of spatial derivatives below are operator norms, which measure the
largest output produced by an input of unit norm. Second derivatives are
bilinear operator norms, with both input vectors restricted to unit norm.
The time interval is $[0,T]$, with $0<T\le1$. Write $\gamma_d=\mathcal N(0,I_d)$.
The class $\mathcal B_K$ and its $C^1_{\rm loc}$ closure $\overline{\mathcal B}_K$
are those of Eq.~\eqref{eq:budget}. The closure includes limits of sequences
whose fields and first derivatives converge uniformly on each compact
space--time set. We include both endpoints of the time interval and use
one-sided time derivatives there.
A limiting field need not have continuous second derivatives.

\subsection{Constants and deterministic estimates}
We first control how far particles move, how the exact flow depends on its
starting point, and how Euler error accumulates. The constants depend only
on the common derivative bound and source moments. They therefore remain
valid when the fitted field changes.
Set
\begin{align}
 C_F&=e^{KT}, & A&=K(1+K), & B&=K(1+K)^2,\\
 J&=B+KA, & C_D&=K+2K^2, & C_E&=AC_F^3,\\
 M_2&=2(1+d), & M_4&=8\{1+d(d+2)\}, & M_3&=M_4^{3/4}.
 \label{eq:base-constants}
\end{align}
These are conservative upper bounds. We use $B$ only for this scalar constant;
model-path bias is denoted $\mathfrak b(u)$.
The variational equations describe the derivatives of the flow with respect
to its starting point. Applying Gronwall's inequality to those equations and
to the trajectories gives, for the flow $F_{s,t}$ of any $u\in\mathcal B_K$,
\begin{align}
 1+|F_{s,t}(x)|&\le e^{K|t-s|}(1+|x|), &
 \|DF_{s,t}\|&\le e^{K|t-s|},\\
 \|D^2F_{s,t}\|&\le K|t-s|e^{2K|t-s|}, &
 |a_t(x)|&\le A(1+|x|),\\
 \|Da_t(x)\|&\le C_D(1+|x|), &
 |j_t(x)|&\le J(1+|x|)^2,
 \label{eq:flow-bounds}
\end{align}
where $a=\partial_tu+Du\,u$ is the acceleration along an exact trajectory
and $j=\partial_ta+Da\,u$ is its time derivative along that trajectory,
sometimes called jerk.
The Jacobian $DF_{s,t}$ records how a small change in the starting point
changes the final position. The Hessian records how this Jacobian changes.
To obtain its bound, differentiate the Jacobian equation and integrate
$\|D^2u\|\|DF\|^2$ using the solution of the associated homogeneous linear
equation. The next identity follows by expanding the derivative along a
trajectory.
\begin{equation}
 j-Du\,a=b:=\partial_{tt}u+2D\partial_tu\,u+D^2u[u,u],
 \qquad |b(t,x)|\le B(1+|x|)^2
 \label{eq:transport-jerk}
\end{equation}
The right side contains exactly the second derivatives controlled by
Eq.~\eqref{eq:budget}.

To compare individual trajectories, start the exact flow and Euler from the
same random point $G\sim\gamma_d$. This pairs their outputs and thus defines
a coupling of the two distributions. The root mean squared displacement
under this pairing gives an upper bound on Wasserstein distance. The
pairing need not be optimal. If
$S_{2,j}=\sum_{k<j}h_k^2$, then
\begin{align}
 1+|X_{t_j}|,\ 1+|\widehat X_j|&\le C_F(1+|G|),\\
 |\widehat X_j-X_{t_j}|&\le C_E S_{2,j}(1+|G|).
 \label{eq:strong}
\end{align}
Here the particle defect means the difference between one Euler update and
one exact update from the same position. It is at most
$Ah^2e^{Kh}(1+|x|)/2$ and the Euler map has Lipschitz constant $1+Kh$.
The Lipschitz bound controls the propagation of errors from earlier steps.
Combining the one-step defects with these propagation factors gives the
displayed $C_E$ bound.
The moments satisfy $\E(1+|G|)^2\le M_2$,
$\E(1+|G|)^4\le M_4$, and $\E(1+|G|)^3\le M_3$.
These are moments of the Gaussian starting point, not moments of a target law.

The next lemma allows us to take limits of fitted fields. It also ensures
that their exact and numerical distributions converge to those of the
limiting field.

\begin{lemma}[Compact closure and continuous laws]
\label{lem:compact}
The class $\overline{\mathcal B}_K$ is sequentially compact in $C^1_{\rm loc}$.
Its elements are $C^1$, with locally Lipschitz first derivatives, globally
$K$-Lipschitz spatial velocities, and the growth bounds in
Eq.~\eqref{eq:budget}. They have complete $C^1$ bi-Lipschitz flows on the time
strip. If $u_n\to u$ in $C^1_{\rm loc}$, then the flows and their spatial
Jacobians converge uniformly on compact cylinders, and
\[
 \sup_{0\le t\le T}\|F^{u_n}_{0,t}(G)-F^u_{0,t}(G)\|_{L^2}\longrightarrow0.
\]
The same $L^2$ continuity holds for the Euler law on any fixed finite grid.
\end{lemma}
\begin{proof}
We first prove convergence on bounded space--time regions. We then control
the contributions from starting points outside those regions. A compact cylinder here
is a compact time interval times a bounded closed spatial set.
On each compact cylinder the first derivatives are bounded and
uniformly Lipschitz, by the three second-derivative bounds in
Eq.~\eqref{eq:budget}. The function values are bounded there by the growth
bound. Arzel\`a--Ascoli, a diagonal subsequence, and the fundamental theorem of
calculus identify a $C^1$ limit with locally Lipschitz first derivatives.
The local topology can be described by a metric. Taking the closure therefore
preserves this compactness property. The common linear-growth and spatial
Lipschitz bounds give flows throughout the interval, including backward
time. Gronwall's inequality on a common compact set of trajectories gives
flow convergence. The Jacobian equation has coefficients $Du_n$ that converge
locally uniformly, so the same argument gives convergence of the Jacobians.

For $|G|\le R$, the convergence is uniform in time. On the complement,
the square of the difference is at most $4C_F^2(1+|G|)^2$, an integrable
Gaussian tail independent of $n$. Taking $R\to\infty$ proves the asserted
$L^2$ convergence. On a fixed grid, Euler applies only finitely many update
maps. These maps converge locally uniformly and satisfy the same growth
bound, so the same argument applies to their composition. Bounds
\eqref{eq:flow-bounds} involving first derivatives pass directly to the limit.
Second-derivative bounds hold almost everywhere. The Taylor and weak-error
estimates below pass from smooth approximations because their remainders
have common bounds. This argument does not require $C^2$ convergence.
\end{proof}

We also need to pass from accuracy at grid nodes to accuracy at every time.
For a prescribed $W_2$-continuous path $p$ starting at $\gamma_d$, write
\[
 \mathfrak b(u)=\sup_{t\in[0,T]}W_2(\rho_t^u,p_t),\qquad
 \omega_p(r)=\sup_{|t-s|\le r}W_2(p_t,p_s).
\]
The first quantity is the largest continuous-model bias along the path.
The second bounds how much the prescribed distribution can change over a
short time interval. Compactness and the preceding lemma make $\mathfrak b$
continuous on the closure. With $E_c$ defined at the coarse nodes,
Eq.~\eqref{eq:strong} gives
\begin{equation}
 \mathfrak b(u)\le E_c+C_E\sqrt{M_2}\,TH+KC_F\sqrt{M_2}\,H+\omega_p(H).
 \label{eq:bias-general}
\end{equation}
The second term bounds integration error at nodes, and the final two terms
compare a time with the preceding node. For the independent interpolation
\eqref{eq:independent}, $\omega_p(H)\le\sqrt{d+\nu}\,H$, where
$\nu=\tr S$. Consequently
\begin{equation}
 \mathfrak b(u)\le E_c+C_{\rm path}H,\qquad
 C_{\rm path}=C_E\sqrt{M_2}\,T+KC_F\sqrt{M_2}+\sqrt{d+\nu}.
 \label{eq:bias-from-nodes}
\end{equation}
This bound needs no exact marginal preservation by $u$.

\section{Signed recursive errors and the two-grid contrast}
\label{app:contrast}

A test average can change in either direction. We keep its sign so we
can tell when changes from different steps add or cancel. To use the
same measurement at different times, trace a point backward through the
exact flow before applying the starting test. The value is constant
along each exact trajectory, so the exact-model average is shared by
both runs. Subtracting their errors cancels that average.\AIassistance

Let $\phi\in C^{1,1}(\R^d)$ with
$\|\nabla\phi\|_\infty\le1$ and $\operatorname{Lip}(\nabla\phi)\le1$.
Set
\begin{equation}
 \phi_t=\phi\circ F_{t,0},\qquad
 I_u(t)=\E[\nabla\phi_t(X_t)\cdot a_t(X_t)].
 \label{eq:I-def}
\end{equation}
The exact-flow expectation $\E\phi_t(X_t)=\E\phi(G)$ is constant.
The quantity $I_u$ measures the effect of acceleration on this time-dependent
test. The next constants control its variation and the remainders in the
expectation error. Such an error is called a weak error because it compares
averages of tests, rather than the positions of individual particles. Define
\begin{align}
 H_\phi&=(1+KT)C_F^3,\\
 C_0&=JC_F^3/6+KAC_F^2/2+TH_\phi A^2C_F^2/8,\\
 C_b&=H_\phi A+C_FC_D,\\
 C_w&=C_0C_F^2M_2+\tfrac12 C_bC_FC_EM_2T,
 & M_w&=BC_F^3\sqrt{M_4}.
 \label{eq:weak-constants}
\end{align}

\begin{lemma}[Uniform weak error, including subintervals]
\label{lem:weak}
For every pair of nodes $i<j$, write
$D_k=\E\phi_{t_k}(\widehat X_k)-\E\phi(G)$. Then
\begin{equation}
 D_j-D_i=-\tfrac12\sum_{k=i}^{j-1}h_k^2I_u(t_k)+R_{ij},\qquad
 |R_{ij}|\le C_wH\sum_{k=i}^{j-1}h_k^2.
 \label{eq:weak-interval}
\end{equation}
Moreover $\operatorname{Lip}(\phi_t)\le C_F$ and
$\operatorname{Lip}(I_u)\le M_w$, uniformly over fields in
$\overline{\mathcal B}_K$.
\end{lemma}
\begin{proof}
We first compare a numerical update and an exact update from the same point.
We then include the error already present in the numerical starting point.
This second comparison is needed because Euler does not restart from the
exact distribution at each step.
First take a smooth field. For $Y=F_{t,t+h}(x)$ and $E=x+hu(t,x)$,
trajectory Taylor expansion gives
\[
 E-Y=-\tfrac12h^2a_t(x)+R,\qquad
 |R|\le JC_F^2h^3(1+|x|)^2/6,
\]
and the lower-order estimate is $|E-Y|\le A C_F h^2(1+|x|)/2$.
The transported test has Hessian bound $H_\phi$, and
\[
 \nabla\phi_{t+h}(Y)=DF_{t,t+h}(x)^{-\top}\nabla\phi_t(x).
\]
The difference from $\nabla\phi_t(x)$ is at most $KhC_F$.
The $C^{1,1}$ Taylor inequality therefore proves
\[
 \phi_{t+h}(E)-\phi_t(x)
 =-\tfrac12h^2\nabla\phi_t(x)\cdot a_t(x)+R_0,
 \quad |R_0|\le C_0h^3(1+|x|)^2.
\]
The three terms in $C_0$ bound the three remainder contributions in this
one-step comparison.
The spatial gradient of $\nabla\phi_t\cdot a_t$ has norm at most
$C_b(1+|x|)$ almost everywhere. Comparing its numerical-state and exact-state
expectations by \eqref{eq:strong} costs at most
$C_bC_FC_EM_2 S_{2,k}\le C_bC_FC_EM_2TH$ at every node, including nodes after $i$.
Now sum over the requested subinterval to obtain \eqref{eq:weak-interval}.
The common $TH$ bound includes particle error accumulated before $i$.
Thus this is an estimate for the actual recursive Euler computation, not
for a computation restarted from exact marginals.

Let $J_t(x)=DF_{0,t}(x)$ and $U_t=J_t^{-1}a_t\circ F_{0,t}$.
Equation~\eqref{eq:transport-jerk} gives
$\partial_tU_t=J_t^{-1}b_t\circ F_{0,t}$ and
$\|\partial_tU_t(G)\|_2\le BC_F^3\sqrt{M_4}$. Since
$I_u(t)=\E[\nabla\phi(G)\cdot U_t(G)]$, this proves the Lipschitz bound.
Smooth approximations of the test preserve its first-derivative and
Lipschitz bounds. To pass to the $C^1$ closure, use
Lemma~\ref{lem:compact} and local convergence of accelerations on bounded
sets. Gaussian moment bounds control the remaining tails. The time-Lipschitz
bound passes as an inequality. It does not assert that a limiting second
derivative is continuous.
\end{proof}

Let $\widetilde X_{2k}$ denote the recursive Euler state on the bisected grid
at the common node $t_k$, for the \emph{same} field. Define
\[
 Z_k=\E\phi_{t_k}(\widehat X_k)-\E\phi_{t_k}(\widetilde X_{2k}).
\]

\begin{lemma}[Two-grid contrast]
\label{lem:contrast}
For any $i<j$,
\begin{equation}
 Z_j-Z_i=-\tfrac14\sum_{k=i}^{j-1}h_k^2 I_u(t_k)+\mathcal R_{ij},\qquad
 |\mathcal R_{ij}|\le C_{\times,w}H\sum_{k=i}^{j-1}h_k^2,
 \quad C_{\times,w}=\tfrac54C_w+M_w/16.
 \label{eq:contrast-full}
\end{equation}
Also $|Z_k|\le C_F(E_c+E_f)$, without assuming that $\rho_t^u=p_t$.
\end{lemma}
\begin{proof}
The fine-grid leading term contributed by one coarse interval is
$-h_k^2\{I_u(t_k)+I_u(t_k+h_k/2)\}/8$. Subtract it from the coarse
term $-h_k^2I_u(t_k)/2$ to obtain
\[
 -h_k^2I_u(t_k)/4+
 h_k^2\{I_u(t_k+h_k/2)-I_u(t_k)\}/8.
\]
The second summand is at most $M_wh_k^3/16$ in absolute value.
The fine-grid remainder has bound $C_w(H/2)(\sum h_k^2/2)$, one quarter
of the coarse bound. This proves the constant in
\eqref{eq:contrast-full}. Finally, couple the two numerical laws and use
$\operatorname{Lip}(\phi_{t_k})\le C_F$, followed by the triangle inequality
through $p_{t_k}$. The prescribed law need not be the exact law of the field.
Comparing one step with two half steps is the classical step-doubling
construction. Here it removes the candidate model's unknown bias and gives
a bound that is uniform over the common regularity class.
\end{proof}

\section{Proof of the two-grid characterization}
\label{app:characterization}

The proof first shows that sufficiently accurate coarse and fine runs must
approach an exact generator whose acceleration preserves the distribution.
It then shows how such a generator permits second-order distribution
accuracy. We keep the compact closure because a convergent sequence of
smooth fields need not have a smooth limit. The proof permits unequal
deterministic step sizes.\AIassistance

\subsection{Necessity: a successful sequence produces an order-elevating generator}
First we identify a limiting exact generator. We then show that a test with
a nonzero average acceleration effect would contradict the assumed accuracy
of the two runs. The contradiction uses a time interval on which that effect
has one sign. This is why accuracy is required along the path, rather than
only at its endpoint.

Suppose there exist $u_n\in\mathcal B_K$, grids with $N_n\to\infty$ and
$H_n\to0$, and $N_n(E_c+E_f)\to0$. Extract $u_n\to u$ in
$C^1_{\rm loc}$ using Lemma~\ref{lem:compact}. By
\eqref{eq:bias-general}, $\mathfrak b(u_n)\to0$; continuity of the
exact laws shows that $u$ generates $p$ on the whole interval.

Assume the acceleration condition fails at some $s\in[0,T]$. Then there
is a compact smooth $\psi$ for which
$\E[\nabla\psi(X_s)\cdot a_s(X_s)]\ne0$.
Set $\phi=\psi\circ F^u_{0,s}$ and rescale by a positive constant so that
its gradient and gradient-Lipschitz bounds are at most one. The flow is
$C^{1,1}$ in space, so this is an admissible compactly supported $C^{1,1}$
test. Its transported version equals a positive multiple of $\psi$ at $s$.
For each $u_n$, transport this \emph{same source test} with that field's
own flow. Lemma~\ref{lem:compact} and the representation
$I_{u_n}(t)=\E[\nabla\phi(G)J_{n,t}^{-1}a^{u_n}_t(X^{u_n}_t)]$
show that $I_{u_n}\to I_u$ uniformly in $t$. Local convergence is uniform
on bounded Gaussian initial states; the tail is bounded by a common
multiple of $1+|G|$.

Thus there is a closed interval of positive length $\ell$, possibly
one-sided at an endpoint, on which all sufficiently large $n$ have
$I_{u_n}$ of one sign and $|I_{u_n}|\ge\eta>0$.
Choose the first and last coarse nodes inside this interval. When $H_n$
is small their separation is at least $\ell/2$. Equation~\eqref{eq:contrast-full}
and $C_{\times,w}H_n\le\eta/8$ yield
\[
 |Z_j-Z_i|\ge\frac\eta8\sum_{k=i}^{j-1}h_k^2
 \ge\frac{\eta\ell^2}{32N_n}.
\]
The last inequality follows from Cauchy--Schwarz. It needs neither a positive
lower bound on individual step sizes nor a limit for their relative sizes.
But the lemma also gives
$|Z_j-Z_i|\le2C_F(E_c+E_f)=o(N_n^{-1})$, a contradiction.
Therefore $\int\nabla\psi\cdot a_t\,dp_t=0$ for every compact smooth
$\psi$ and every $t$. This proves necessity for arbitrary refining grids,
not only uniform or quasi-uniform ones.

\subsection{Sufficiency: locally Lipschitz acceleration is enough here}
Euler omits the leading displacement due to acceleration. To compare
distributions, we can pair an exact particle with a different numerical
particle, provided each population still has the same distribution.
The acceleration flow gives such a pairing. With physical time fixed,
it supplies the leading displacement Euler misses while leaving the
input distribution unchanged. This change of
pairing is a recoupling.

Suppose $u\in\overline{\mathcal B}_K$ generates $p$ and satisfies the
acceleration condition. Fix $t$, and let $\Psi_t^r$ be the autonomous flow
of $a_t$. The acceleration is locally Lipschitz with linear growth, so
this flow is complete and locally bi-Lipschitz. It preserves $p_t$.
For clarity, $p_t$ is the positive continuous density obtained by pushing
Gaussian density through a $C^1$ bi-Lipschitz flow. The weak identity
$\operatorname{div}(p_ta_t)=0$ extends from smooth tests to compactly
supported Lipschitz tests by local mollification. Applied to
$\psi\circ\Psi_t^r$, the characteristic chain rule gives
$\frac{d}{dr}\int\psi(\Psi_t^r(x))p_t(dx)=0$ almost everywhere in $r$.
Preimages of compact sets remain compact on bounded $r$ intervals by
linear growth. The derivative and the local mollification limits are
therefore dominated. Integrating in $r$ proves invariance. Equivalently,
this is uniqueness of the stationary continuity equation under a locally
Lipschitz complete vector field.

For $X\sim p_t$ use the recoupling $Z=\Psi_t^{h^2/2}(X)$.
The new starting point has the same distribution, $Z\sim p_t$.
Euler starting at $Z$ therefore has the required one-step output law, but
it is now paired differently with the exact particle. Let $G_A=e^{AT^2/2}$.
The acceleration estimates imply
\begin{align*}
 |Z-X|&\le (A G_A/2)h^2(1+|X|),\\
 |Z-X-\tfrac12h^2a_t(X)|
 &\le (C_D A G_A^2/8)h^4(1+|X|)^2.
\end{align*}
The second inequality follows by integrating
$|a_t(\Psi_t^rX)-a_t(X)|$ and using the local bound on $Da_t$ along
the orbit. Exact-flow Taylor expansion gives
\[
 F_{t,t+h}(X)=X+hu_t(X)+\tfrac12h^2a_t(X)+R,
 \quad |R|\le JC_F^2h^3(1+|X|)^2/6.
\]
This remains valid in the closure by passage from smooth fields with the
uniform remainder. Consequently
\begin{equation}
 |Z+hu_t(Z)-F_{t,t+h}(X)|
 \le C_{\rm loc}h^3(1+|X|)^2,
 \quad C_{\rm loc}=C_D A G_A^2T/8+KA G_A/2+JC_F^2/6.
 \label{eq:AP-local}
\end{equation}
We can now average the particle bound using the Gaussian fourth moment.
This gives local $W_2$ error at most
$C_{\rm loc}C_F^2\sqrt{M_4}h^3$. Propagating it with Euler's
$(1+Kh)$ Lipschitz stability proves, on every grid,
\begin{equation}
 \max_k W_2(\widehat\rho_k^u,p_{t_k})
 \le C_F^3 C_{\rm loc}\sqrt{M_4}\sum_kh_k^3
 \le C_F^3 C_{\rm loc}\sqrt{M_4}\,TH^2.
 \label{eq:AP-global}
\end{equation}
The acceleration need not be globally Lipschitz: its derivative grows at
most linearly, and the required fourth moment comes from Gaussian
initialization and the common linear-growth bound on the velocity.

If $u$ belongs only to the closure, approximate it by a smooth
$u_N\in\mathcal B_K$ so closely that both finite-grid Euler laws have
additional $W_2$ discrepancy at most $N^{-2}$ at their finitely many common
nodes. Such approximation follows from Lemma~\ref{lem:compact}; it need
not preserve the target path exactly. On the uniform grid and its
bisection, \eqref{eq:AP-global} proves $\mathcal J_N(u_N)=O(N^{-2})$.
The same approximant is used for the two runs. This proves sufficiency
without asserting that a $C^1$ limit is itself a smooth member of the class.

\subsection{Uniform lower bound and the three optimized regimes}
The same necessity argument gives a lower bound for all admissible fields
on sufficiently fine grids. If no such common bound existed, we could
choose a sequence with progressively smaller error. That sequence would
produce exactly the generator that is assumed not to exist.

If there were no constants $c,H_0>0$ in
\eqref{eq:general-uniform-barrier}, for each integer $m$ we could select a
field and a grid with $H_m\le1/m$ and
$\mathcal J_{N_m}<1/(mN_m)$. Necessarily $N_m\ge T/H_m\to\infty$.
The necessity argument would produce an order-elevating generator in
the closure, a contradiction. Thus the lower bound is uniform over all
fields, all sufficiently fine deterministic grids, and all numbers of
steps, not just over a single refinement sequence.

For uniform grids define $q_N=\inf_{u\in\mathcal B_K}\mathcal J_N(u)$.
For each fixed $N$ this infimum equals a minimum over the compact closure,
since Euler laws are continuous there. Put
$b_*=\min_{u\in\overline{\mathcal B}_K}\mathfrak b(u)$.
Uniformly over that closure,
\[
 |E_c-\mathfrak b(u)|+|E_f-\mathfrak b(u)|
 \le C h+2\omega_p(h),\qquad h=T/N,
\]
for a class-dependent constant $C$. This follows from strong Euler error
and the difference between a continuous-time supremum and the maximum
over common coarse nodes. Therefore $q_N\to2b_*$.
If no exact generator exists, compactness gives $b_*>0$.
If an exact generator exists but no order-elevating generator does, the
uniform lower bound gives $q_N\ge c/N$, and approximation of an exact
generator, together with strong Euler error, gives $q_N\le C/N$.
If an order-elevating generator exists, the previous subsection gives
$q_N=O(N^{-2})$. This is a statement about the optimized two-grid error;
individual nonoptimal sequences can have other rates.

\section{Gaussian-source topology and the odd-dimensional application}
\label{app:source}

We now specify a path for which the generator required by
Theorem~\ref{thm:characterization} cannot exist in odd dimension.
Only the first two time derivatives of the path at its Gaussian starting
distribution are needed. This pair of derivatives is called the source
two-jet. They constrain every exact generator, including nonlinear ones.
We can therefore prove nonexistence without first listing all possible
generators. The target need not be Gaussian, radial, or unimodal.\AIassistance

\begin{lemma}[Source two-jet identity]
\label{lem:twojet}
Suppose a $u\in\overline{\mathcal B}_K$ generates
$p_t=\Law((1-t)G+tY)$ near zero, where $Y$ is centered, independent of
$G$, and has covariance $S$. Let $r(x)=u(0,x)+x$.
Then $\operatorname{div}(\gamma_dr)=0$. For a unit vector $e$, define
$Q_e(y)=\E[(e\cdot r(G))^2\mid e\cdot G=y]$ and let $\varphi$ be the
standard one-dimensional Gaussian density. For a scalar test $\psi$ with
bounded first and second derivatives and at most linear growth,
\begin{equation}
 \E[\psi'(e\cdot G)e\cdot a_0(G)]
 =\int\psi''(y)\{e^\top Se-Q_e(y)\}\varphi(y)\,dy.
 \label{eq:twojet-proof}
\end{equation}
The equality also holds for $C^{1,1}$ tests with almost-everywhere second
derivatives, by approximation.
\end{lemma}
\begin{proof}
An exact generator and the prescribed interpolation give the same average
of each test wherever their distributions agree. Their first and second
time derivatives must therefore agree at zero. For the first derivative, a compact smooth test
along the prescribed straight interpolation gives
$-\E\nabla\Phi(G)\cdot G$. The derivative along the flow gives
$\E\nabla\Phi(G)\cdot u_0(G)$.
Equality says that the correction leaves the Gaussian source unchanged,
which is Gaussian solenoidality. Testing along $e$ shows
that the distributional derivative of
$\varphi(y)\E[e\cdot r(G)\mid e\cdot G=y]$ vanishes.
This function is integrable, hence the constant is zero.
Lipschitz growth of $r$ and Gaussian conditional expectation give continuous
versions of the conditional first and second moments.

For the second derivative, the value along the straight interpolation is
$\E\psi''(e\cdot G)(e\cdot(Y-G))^2$, whose conditional squared factor is
$y^2+e^\top Se$. Along the flow the second derivative is
\[
 \E\big[\psi''(e\cdot G)(e\cdot u_0(G))^2
       +\psi'(e\cdot G)e\cdot a_0(G)\big].
\]
The first conditional squared factor equals $y^2+Q_e(y)$ because the
conditional mean of $e\cdot r$ is zero. Subtracting the two second-derivative
expressions proves the identity. A $C^1$ field is enough to differentiate
the trajectory twice. Linear growth gives the integrable bounds needed to
differentiate its expectation. Only the second moment of $Y$ is used.
For noncompact scalar tests, first apply spatial cutoffs and then remove
them using Gaussian integrability of $|r|^2$, linear growth of $a_0$, and
the bounded test derivatives. Smooth approximation gives the $C^{1,1}$
case without boundary atoms, since the scalar Gaussian has a density.
\end{proof}

\begin{lemma}[Uniform positive strip test function in odd dimension]
\label{lem:pole}
Assume $d$ is odd and $\lambda=\lambda_{\min}(S)>0$. Define
\begin{equation}
 A_d=d+\sqrt{d^2+d-2},\qquad
 R_K=\max\{1,2A_d(K+1)/\sqrt\lambda\},\qquad
 \delta_K=\frac{\lambda}{2\sqrt{2\pi}}e^{-(R_K+1)^2/2}.
 \label{eq:pole-constants}
\end{equation}
Every exact generator $u\in\overline{\mathcal B}_K$ has a unit direction $e$
and a convex strip test $\phi_e(x)=\psi(e\cdot x)$, with gradient and
gradient-Lipschitz bounds one, such that
$\E[\nabla\phi_e(G)\cdot a_0(G)]\ge\delta_K>0$.
\end{lemma}
\begin{proof}
The identity compares target variance with the mean squared correction in a
chosen direction. We seek a region where the latter is too small to cancel
the former. A topological argument selects a point on a large sphere.
Lipschitz continuity controls the correction near that point. A convex test
then detects the acceleration that remains.

Put $L_r=\operatorname{Lip}(r)\le K+1$. The continuous weighted-divergence
identity is pointwise and gives
$\operatorname{div}r(x)=x\cdot r(x)$, so $|x\cdot r(x)|\le dL_r$.
Project the correction onto directions tangent to the sphere of radius
$R_K$. This gives a continuous tangent vector field on a sphere of even
dimension $d-1$. The hairy-ball theorem says that this field must vanish
somewhere~\citep{hatcher2002algebraic}. In $d=1$ the tangent space is
zero-dimensional, so the step is immediate. Thus $r(R_Ke)$ is parallel to
$e$ for some unit $e$. We have not assumed that $r$ is tangent elsewhere.

For $y\in[R_K,R_K+1]$ write $G=ye+Z_\perp$ under the scalar conditioning,
where $Z_\perp\sim\mathcal N(0,I_{d-1})$. Since
$Z_\perp\cdot r(R_Ke)=0$, the divergence identity and Lipschitz bound give
\[
 |e\cdot r(ye+Z_\perp)|\le\frac{L_r}{R_K}
 \left[d+|Z_\perp|\sqrt{(y-R_K)^2+|Z_\perp|^2}\right].
\]
Use $\E|Z_\perp|^2=d-1$ and
$\E|Z_\perp|^4=(d-1)(d+1)$ in Minkowski's inequality to obtain
$\sqrt{Q_e(y)}\le A_dL_r/R_K\le\sqrt\lambda/2$ throughout the strip.
Choose $\psi'$ equal to zero below $R_K$, increasing with slope one
on the strip, and equal to one above it. Equation~\eqref{eq:twojet-proof}
then gives
\[
 \E\nabla\phi_e(G)\cdot a_0(G)
 \ge\frac{3\lambda}{4}\Pr\{R_K\le N(0,1)\le R_K+1\}
 \ge\frac{3\lambda}{4\sqrt{2\pi}}e^{-(R_K+1)^2/2}\ge\delta_K.
\]
The test average is positive, but its lower bound includes a Gaussian tail
probability. It can therefore be extremely small for large derivative
bounds. The argument establishes positivity, not a useful numerical
forecast of the error.
\end{proof}

\begin{proof}[Proof of Theorem~\ref{thm:odd-transfer}]
An order-elevating generator would give zero average acceleration effect
for every compact smooth gradient test. Cutting off the strip test on
larger and larger sets and controlling the Gaussian tail gives the same
zero value for the strip test in Lemma~\ref{lem:pole}. This contradicts
its positive average.
The general uniform lower bound in Theorem~\ref{thm:characterization}
therefore applies to every field, including fields whose exact law differs
from the prescribed interpolation.

One constant in this lower-bound argument remains nonconstructive: we prove
that it exists without obtaining a formula for it. Compactness yields
$\eta_K>0$ such that $\mathfrak b(u)\le\eta_K$ implies a strip test function
with source pairing at least $\delta_K/2$. Otherwise a sequence of biases
converging to zero would have a $C^1_{\rm loc}$ limit generating $p$;
choose that limit's direction and pass its acceleration pairing to the
sequence to obtain a contradiction. Fix
\[
 \tau=\min\{T,\delta_K/(4M_w)\},\quad
 H_0=\min\{\eta_K/(2C_{\rm path}),\tau/2,
                  \delta_K/(32C_{\times,w})\}.
\]
If $E_c\ge\eta_K/2$, the lower bound is immediate with $c\le\eta_K/2$.
Otherwise \eqref{eq:bias-from-nodes} gives $\mathfrak b(u)\le\eta_K$.
The transported test function has $I_u(t)\ge\delta_K/4$ on $[0,\tau]$.
Let $t_j$ be the last node before $\tau$. The prefix version of
\eqref{eq:contrast-full}, with $Z_0=0$, gives
\[
 C_F(E_c+E_f)\ge|Z_j|\ge\frac{\delta_K}{32}\sum_{k<j}h_k^2
 \ge\frac{\delta_K\tau^2}{128N}.
\]
Thus an admissible $c$ is
$\min\{\eta_K/2,\delta_K\tau^2/(128C_F)\}$.
The strip location $R_K$ and lower bound $\delta_K$ have explicit formulas,
but $\eta_K$ does not. The proof therefore does not give a computable
threshold for the general odd-dimensional result.
\end{proof}

\subsection{A non-Gaussian target with an admissible exact generator}
The lower bound is not only about paths that no regular field can generate.
The following mixture has a regular exact generator. In odd dimension, the
obstruction still prevents its acceleration from satisfying the required
distribution-preserving condition.

For $Y=\mu\xi e_1+\sigma Z$, with independent symmetric sign $\xi$,
standard $d$-dimensional Gaussian $Z$, and $\sigma>0$, let
$q(t)=(1-t)^2+\sigma^2t^2$. The canonical conditional velocity is
\begin{equation}
 v_t(x)=\frac{(1+\sigma^2)t-1}{q(t)}x+
 \frac{\mu(1-t)}{q(t)}
 \tanh\left(\frac{\mu t x_1}{q(t)}\right)e_1.
 \label{eq:mixture-generator}
\end{equation}
This velocity generates the independent path: it averages $Y-G$ conditional
on the interpolated position. The function $q$ stays
strictly positive on a neighborhood of $[0,1]$, so all scalar coefficient
derivatives needed here are bounded. Write the nonlinear term
$c(t)\tanh(b(t)x_1)$. Its time derivative grows at most linearly in $x$;
its second time derivative at most quadratically. Spatial derivatives
through order two are uniformly bounded. For the mixed derivative,
terms such as $b(t)x_1\tanh''(b(t)x_1)$ are uniformly bounded, since
$z\tanh''z$ is bounded on the real line. All remaining terms are bounded
scalar coefficients times bounded hyperbolic derivatives. Hence this
field belongs to some finite $\mathcal B_K$, including for separated
multimodal targets.
Thus the class of exact generators can be nonempty for a non-Gaussian
target. Positive covariance alone still does not make every
$\mathcal P_2$ target reachable by a globally Lipschitz flow.

\section{An explicit initial-risk bound without exact marginal preservation}
\label{app:risk}

Throughout this section the target path is \eqref{eq:independent}, but
$u$ is any field in $\overline{\mathcal B}_K$. Put $\nu=\tr S$ and
$\Delta_0(u)=\E|u(0,G)+G|^2$. Independence gives the elementary risk identity
\begin{equation}
 \mathcal L_0(u)=\E|u(0,G)-(Y-G)|^2=\nu+\Delta_0(u).
 \label{eq:source-loss-proof}
\end{equation}
Here population loss means an expectation over the source and target
distributions, rather than an average over a finite training batch.
The Bayes velocity minimizes this squared prediction loss. At the source
it is $-x$, with minimum loss $\nu$. The excess risk $\Delta_0$ measures
how much loss the candidate adds above that minimum. The decomposition is
standard. We use it to constrain a candidate through its two numerical runs,
without assuming that its continuous flow preserves the target path.\AIassistance

\subsection{Quadratic transported tests and all constants}
We now test the mean squared position, which is the second moment. This
test is unbounded and its gradient grows linearly. We therefore keep the
moment factors explicitly when repeating the earlier comparison of test
averages. The constants below make the final source-risk inequality
computable.

Use \eqref{eq:base-constants} and set
\begin{align}
 H_q&=(1+KT)C_F^3, & D_q&=3C_F^3\sqrt{M_2}, & M_q&=BC_F^4\sqrt{M_2M_4},\\
 D_0&=AC_F^2M_2, & c_\ell&=\min\{T/D_0,1/(2M_q)\},\\
 C_{q0}&=JC_F^5/6+KAC_F^3/2+TH_qA^2C_F^3/4,\\
 C_{qb}&=H_qA+C_F^2C_D, &
 C_q&=M_3\{C_{q0}C_F^3+\tfrac12 C_{qb}C_F^2C_ET\},\\
 C_{\times,q}&=\tfrac54 C_q+M_q/16, &
 C_H&=\max\{2/c_\ell,16C_{\times,q}\},\\
 C_m&=C_F\sqrt{M_2}+\sqrt d+\sqrt\nu, &
 J_m&=6KAC_F^2M_2+2JC_F^3M_3.
 \label{eq:risk-constants}
\end{align}
Let $\psi_t(x)=|F_{t,0}(x)|^2/2$ and
$I(t)=\E[\nabla\psi_t(X_t)\cdot a_t(X_t)]$.
Then
\begin{align}
 |\nabla\psi_t(x)|&\le C_F^2(1+|x|), &
 \|D^2\psi_t(x)\|&\le H_q(1+|x|),\\
 |I'(t)|&\le M_q, & |I(0)|&\le D_0.
 \label{eq:quadratic-bounds}
\end{align}
The derivative bound follows from \eqref{eq:transport-jerk}. Take the dot
product of transported acceleration with $G$, then use Cauchy--Schwarz and
the Gaussian moment bounds. The same constant also gives a conservative
bound when the test is transported to any other fixed reference time.

The proof of Lemma~\ref{lem:weak}, now with a linearly growing Hessian,
gives a one-step remainder at most $C_{q0}h^3(1+|x|)^3$.
To see the three terms, the trajectory Taylor remainder costs
$JC_F^5h^3(1+|x|)^3/6$, the transport of the gradient costs at most
$KAC_F^3h^3(1+|x|)^2/2$, and the quadratic Taylor remainder costs at most
$TH_qA^2C_F^3h^3(1+|x|)^3/4$.
The scalar measurement $\nabla\psi_t\cdot a_t$ has derivative bounded by
$C_{qb}(1+|x|)^2$. At numerical and exact states its expectation difference
is at most $C_{qb}C_F^2C_EM_3S_{2,k}$.
Summation and $S_{2,k}\le TH$ yield the interval weak remainder
$C_qH\sum h_k^2$.
Consequently the two-grid contrast formula \eqref{eq:contrast-full} holds
with $C_{\times,q}$ in place of $C_{\times,w}$.

Although $\psi_t$ is not globally Lipschitz, for the two numerical laws
$\alpha,\beta$ at the same time it satisfies
\begin{equation}
 \left|\int\psi_t\,d\alpha-\int\psi_t\,d\beta\right|
 \le D_qW_2(\alpha,\beta)\le D_q(E_c+E_f).
 \label{eq:quadratic-comparison}
\end{equation}
Indeed the mean-value inequality bounds the integrand by
$C_F^2(1+|x|+|y|)|x-y|$; both numerical states have root second moment at
most $C_F\sqrt{M_2}$. Cauchy--Schwarz gives the conservative constant $D_q$.
The target is used only in the Wasserstein triangle inequality, so no
fourth moment of $Y$ is required here.

Put $a_*=|I(0)|=|\E G\cdot a_0(G)|$. If $a_*>0$, on
$[0,\ell]$ with $\ell=\min\{T,a_* /(2M_q)\}\ge c_\ell a_*$ the
pairing has a common sign and magnitude at least $a_*/2$.
If $H\le\ell/2$ and $H\le a_* /(16C_{\times,q})$, choose the last node
before $\ell$. The contrast and \eqref{eq:quadratic-comparison} give
\[
 D_q(E_c+E_f)\ge\frac{a_*}{16}\sum_{k<j}h_k^2.
\]
Cauchy--Schwarz gives
\begin{equation}
 a_*\le\left\{\frac{64D_qN(E_c+E_f)}{c_\ell^2}\right\}^{1/3}.
 \label{eq:cubic-observe}
\end{equation}
On the uniform grid, $\sum_{k<j}h_k^2=(T/N)t_j\ge T\ell/(2N)$, giving
\begin{equation}
 a_*\le\left\{\frac{32D_qN(E_c+E_f)}{Tc_\ell}\right\}^{1/2}.
 \label{eq:square-observe}
\end{equation}
If either small-mesh condition fails, then $a_*\le C_HH$ instead.
Adding this term makes each displayed bound valid without a mesh
qualification. If $a_*=0$, the same statements hold trivially.

\subsection{Recovering the first two source-time moment derivatives}
To use the source-risk identity, we need the first two time derivatives of
the model's second moment at time zero. Accuracy along the path controls
the moment itself. A common bound on its third derivative then controls
these two source-time derivatives. We obtain explicit bounds by using
forward finite differences, which estimate derivatives from nearby times.

Let $m_u(t)=\E|X_t^u|^2$,
$m_p(t)=d(1-t)^2+\nu t^2$, and $f=m_u-m_p$.
Then $f(0)=0$ and
\begin{equation}
 \|f\|_\infty\le C_m\mathfrak b(u),\qquad
 |f'''(t)|\le J_m\quad\hbox{almost everywhere}.
 \label{eq:moment-bias}
\end{equation}
The square roots of second moments differ by at most the Wasserstein
distance. This inequality gives the first bound. The second follows from
$m_u'''=6\E[u_t\cdot a_t]+2\E[X_t\cdot j_t]$ and
\eqref{eq:flow-bounds}. For a closure field the inequality follows by
passing the uniform finite-difference estimates of smooth approximants;
classical continuous third derivatives are not being assumed.

For $0<s\le T/2$, forward differences and the Lipschitz bound on $f''$
give
\begin{align}
 |f'(0)|&\le\frac{5C_m\mathfrak b(u)}{2s}+J_ms^2,\\
 |f''(0)|&\le\frac{3C_m\mathfrak b(u)}{s^2}+J_ms.
 \label{eq:jet-estimates}
\end{align}
The first formula uses $(4f(s)-f(2s))/(2s)$; the second uses
$(f(2s)-2f(s)+f(0))/s^2$.
For the latter, the double-integral kernel has average evaluation time
$s$, so a $J_m$-Lipschitz second derivative incurs at most $J_ms$.
Since $s\le1$,
\begin{equation}
 |f'(0)+f''(0)/2|\le\omega_K(\mathfrak b(u)),\quad
 \omega_K(b)=\inf_{0<s\le T/2}\left\{\frac{4C_mb}{s^2}+\frac32J_ms\right\}.
 \label{eq:omega-proof}
\end{equation}
The infimum chooses the finite-difference spacing that gives the smallest
bound. Here $\omega_K(0)=0$, and for $b>0$ the minimizer is
$s=\min\{T/2,(16C_mb/(3J_m))^{1/3}\}$. This explicit modulus is
nondecreasing and $O(b^{1/3})$ as $b\downarrow0$.

The exact source identity, valid for every field, is
\begin{equation}
 \E[G\cdot a_0(G)]
 =\nu-\Delta_0(u)+f'(0)+\tfrac12f''(0).
 \label{eq:approximate-risk-identity}
\end{equation}
To verify it, note that
$m_u'(0)=2\E[G\cdot u_0(G)]$ and
$\E|u_0(G)|^2=d-f'(0)+\Delta_0(u)$.
Insert these into
$m_u''(0)=2\E|u_0(G)|^2+2\E[G\cdot a_0(G)]$ and use
$m_p''/2=d+\nu$. This calculation does not require Gaussian
solenoidality or exact path preservation.

\begin{proof}[Proof of Theorem~\ref{thm:source-risk}]
Combine \eqref{eq:approximate-risk-identity}, \eqref{eq:omega-proof},
\eqref{eq:bias-from-nodes}, and either \eqref{eq:cubic-observe} or
\eqref{eq:square-observe}, including the $C_HH$ alternative.
This gives precisely Eqs.~\eqref{eq:risk-main}--\eqref{eq:risk-uniform}.
If $H_N\to0$ and $E_c+E_f=o(N^{-1})$, every term tends to zero,
so $\Delta_0(u_N)\to\nu$ and $\mathcal L_0(u_N)\to2\nu$.
For a fixed gap $|\Delta_0-\nu|\ge g>0$, an entirely explicit choice is
\begin{align}
 s_g&=\min\{T/2,g/(12J_m)\}, & b_g&=g s_g^2/(32C_m),\\
 H_g&=\min\{T,g/(4C_H),b_g/(2C_{\rm path})\}, &
 c_g&=\min\{b_g/2,c_\ell^2g^3/(4096D_q)\}.
 \label{eq:explicit-gap-barrier}
\end{align}
If $H\le H_g$ and $\mathcal J_N<c_g/N$, then
$E_c+C_{\rm path}H\le b_g$. Evaluating the infimum in $\omega_K$
at $s_g$ gives at most $g/4$, the mesh term is at most $g/4$, and
the cubic term is at most $g/4$. The asserted gap would therefore be
at most $3g/4$, a contradiction. Thus the risk-gap barrier is uniform
over all sufficiently fine deterministic grids, with no lower step
bound or condition $H=O(N^{-1})$.
Although every constant has a formula, the bounds may be far too
conservative to give a useful numerical certificate.
\end{proof}

\section{Exact examples, loss of a fit after refinement, and necessary restrictions}
\label{app:controls}

\subsection{An exact fit on one Euler grid can fail after refinement}
The next example separates two questions: does Euler follow the fitted
velocity accurately, and does its output match the target? We choose a
velocity for which Euler matches the target exactly at one step size.
The exact continuous flow of that velocity does not match the target.
Using smaller steps therefore exposes the model error that the coarse
numerical error had canceled.\AIassistance

Let $g(t)=\sqrt{(1-t)^2+\sigma^2t^2}$, $\sigma>0$, and
$u_h(t,x)=c_h(t)x$, where
$c_h(t)=(g(t+h)-g(t))/(h g(t))$. On a fixed neighborhood of $[0,1]$,
$g$ is positive, so this family has a common derivative bound as $h\downarrow0$.
The exact algebraic identity
\[
 1+hc_h(t)=g(t+h)/g(t)
\]
shows why the coarse computation is exact at every target node. The scale
factors multiply, and each denominator cancels the preceding numerator.
This gives every target node marginal exactly in every dimension.
The exact scale of the continuous model is
$q_h(t)=\exp\{\int_0^tc_h(s)ds\}$ and Taylor expansion gives
\begin{equation}
 q_h(t)=g(t)\exp\left\{\frac h2\int_0^t\frac{g''(s)}{g(s)}ds+O(h^2)\right\}.
 \label{eq:secant-model}
\end{equation}
Now keep this \emph{same} $u_h$ and use Euler step $h/2$. For a smooth scalar
coefficient $c$, the logarithm of the Euler scale on step $k$ is
\[
 \int_0^t c(s)ds-\frac k2\int_0^t\{c'(s)+c(s)^2\}ds+O(k^2).
\]
To obtain this expression, approximate the integral by its left Riemann
sum and expand $\log(1+kc)$. The coefficient family and the derivatives
used in this calculation have common bounds. The remainder bounds are
therefore uniform in $h$.
With $c_h=g'/g+(h/2)g''/g+O(h^2)$ and $k=h/2$, the refined scale satisfies
\begin{equation}
 \widehat g_{h/2;u_h}(t)
 =g(t)\exp\left\{\frac h4\int_0^t\frac{g''(s)}{g(s)}ds+O(h^2)\right\}
 \label{eq:secant-transfer}
\end{equation}
at common nodes. Since $g''/g=\sigma^2/g^4>0$, the refined target error
is $\Theta(h)$ even though the fitted coarse error is exactly zero.
The same computation for step $h/m$ has coefficient
$(h/2)(1-1/m)$, for any fixed integer $m\ge2$.
Thus a velocity chosen for one grid can attain the first-order lower
bound. This calculation explains the cancellation, but it does not give
a general learning algorithm. If a network receives the requested step
size as an input, changing that input can change the evaluated field.
Such a comparison falls outside the same-field requirement here.

For the convex blend
$c_{h,\theta}=(1-\theta)g'/g+\theta c_h$, the leading logarithmic target
errors on steps $h$ and $h/2$ have coefficients
$(\theta-1)h/2$ and $(\theta-1/2)h/2$ respectively.
Their largest absolute coefficient is minimized at $\theta=3/4$, with
strictly positive value $h/8$. Fitting both resolutions within this scalar
family therefore reduces the error constant while leaving the order
unchanged. The numerical one-dimensional bounded optimization checks
this fact for the stated family, not for all neural fields.

\subsection{Even-dimensional exactness and insufficiency of scalar risk}
The restriction to odd dimension matters. In even dimension, a rotation
can change which source and target points are paired while preserving
the Gaussian distribution at every time. The resulting trajectories can
be straight, so Euler can be exact. The following construction provides
this contrasting case.

For even $d$ choose $J^\top=-J$, $J^2=-I$, and set
$A_t=(1-t)I+\sigma tJ$. The process $A_tG$ has covariance $g(t)^2I$,
with straight trajectories. Its velocity
$u_t=(\sigma J-I)A_t^{-1}x$ has zero pointwise acceleration, belongs to
a class $\mathcal B_K$ with finite $K$, and Euler is exact on every grid.
This skew-coupling construction is known and is credited to the existing
Gaussian geometry rather than counted as a new theorem
\citep{obzherin2026density,obzherin2026finite}.
At the source, $u_0+\operatorname{id}=\sigma Jx$, so
$\Delta_0=d\sigma^2=\nu$ and the doubled-risk necessity is attained.

Matching the required source loss alone is not sufficient to make
acceleration distribution-preserving. In $d=3$, take $S=I_3$,
$g^2=(1-t)^2+t^2$, and
\[
 u_t(x)=\left\{(g'/g)I+\sqrt{3/2}\,g^{-2}J_{12}\right\}x,
\]
where $J_{12}$ is the unit rotation on coordinates $1,2$ and is zero
on coordinate $3$. This is an exact generator. Its source excess is
$3=\nu$, but $a_0(x)=\operatorname{diag}(-1/2,-1/2,1)x$.
The quadratic test $x_3^2/2$ detects a positive average acceleration
effect, so the acceleration condition fails. Thus neither the loss equality
nor small target errors at one or two finite step counts establish the
existence of an order-elevating generator.

\subsection{Why endpoint-only tests are excluded}
An acceleration contribution can be positive during one part of the
interval and negative during another. These contributions can cancel at
the final time while errors remain at intermediate times. The following
scalar example shows why testing only the final distribution cannot
replace testing the whole path.

An explicit scalar example is $X_t=f(t)G$, with
$u(t,x)=(f'(t)/f(t))x$ and $f>0$. The leading terminal Euler coefficient
is proportional to $\int_0^T f''(s)/f(s)\,ds$; it may cancel while the
prefix integrals do not. For example, on $[0,1]$ take
$f(t)=\exp\{-5t^3/3\}$. Then $u(t,x)=-5t^2x$ and
$\int_0^1 f''/f=-5+25/5=0$, although $f''/f$ is not identically
zero. Thus terminal Euler error is second-order for this example while
whole-path error is first-order. The necessity proof compares two times
surrounding an interval on which the average acceleration effect is
nonzero and keeps one sign. Final-time data alone cannot provide those
two-time comparisons. Existing finite-step
Gaussian terminal-cancellation constructions likewise leave intermediate
marginals unconstrained~\citep{obzherin2026finite}. None of the main
results rules out accurate endpoint-only generators, other numerical
methods, state augmentation, or step-conditioned average-velocity models.
